\chapter{Back Pain During Pregnancy}

\section*{Definition}
Back pain during pregnancy is pain in the lower, middle, or upper back occurring during gestation. Mechanical strain is common, but severe, new, radiating, fever-associated, or neurologic pain may signal infection, spinal disease, kidney disease, preterm labor, or another emergency.

\section*{When to See a Doctor}
Seek urgent care for back pain with fever, flank pain, painful urination, vaginal bleeding, fluid leakage, regular contractions, abdominal pain, reduced fetal movement, leg weakness, saddle numbness, or loss of bladder or bowel control. Sudden severe pain after trauma requires assessment.

\section*{How It Develops}
Pregnancy changes posture, shifts the center of gravity, loosens ligaments, and increases load on the spine and pelvic joints. Muscle fatigue and sacroiliac dysfunction can cause pain. Kidney infection, nerve compression, uterine contractions, and rare spinal or vascular conditions can produce more serious pain.

\section*{Causes}
\begin{itemize}
  \item Postural strain, pelvic-girdle pain, sacroiliac dysfunction, or muscle spasm
  \item Sciatica or disk-related nerve-root irritation
  \item Urinary tract infection, kidney infection, or kidney stone
  \item Preterm labor, placental abruption, or other obstetric complication
  \item Fracture, inflammatory disease, or spinal cord compression
\end{itemize}

\section*{Related Symptoms}
\begin{itemize}
  \item Pelvic pressure, contractions, abdominal pain, or vaginal bleeding
  \item Fever, chills, urinary symptoms, nausea, or flank pain
  \item Numbness, weakness, shooting leg pain, or difficulty walking
\end{itemize}

\section*{Medical Evaluation}
\begin{itemize}
  \item History includes gestational age, pain pattern, trauma, fever, urinary symptoms, neurologic symptoms, contractions, bleeding, and fetal movement.
  \item Examination assesses spine, gait, neurologic function, abdomen, costovertebral angle, uterine activity, and fetal status.
  \item Urinalysis, culture, blood tests, fetal monitoring, ultrasound, or MRI is selected according to findings; necessary imaging should not be withheld.
\end{itemize}

\section*{Treatment Options}
Treatment may include physical therapy, activity modification, pregnancy-safe analgesia, antibiotics, treatment of renal disease, or urgent obstetric or neurologic care. Persistent disabling pain warrants reassessment rather than prolonged self-treatment.

\section*{Self-Care and Home Management}
For mild mechanical pain, use posture support, gentle activity, pregnancy-adapted stretching, and heat applied safely to the back. Avoid bed rest, heavy lifting, and unapproved medicines. Ask the prenatal clinician before using analgesics, braces, or exercises.

\section*{References}
\begin{thebibliography}{99}
\bibitem{back_pain_pregnancy:ref1} American College of Obstetricians and Gynecologists. Back Pain During Pregnancy. ACOG patient and clinical guidance, accessed 2025.
\bibitem{back_pain_pregnancy:ref2} American College of Obstetricians and Gynecologists. Diagnostic Imaging During Pregnancy and Lactation. Committee Opinion No. 723, reaffirmed 2021.
\end{thebibliography}
